% ================================================================
\section{The Time-Domain Bridge: Witness Waves and Ramanujan Invariance}
\label{sec:time-domain}
% ================================================================

The Time-Domain Bridge (\lean{Physics/Bridges/TimeDomainBridge.lean},
1,444 lines, zero sorry, one analysis axiom) establishes that the Gram
quadratic form---the central object measuring the vacuum energy---\emph{is}
a time-domain integral of a physical wave. This transforms the abstract
linear algebra of $v^T G v$ into a concrete question about wave interference.

\subsection{The Witness Wave}

Define the \emph{witness wave}:
\begin{equation}
  V(t) = \sum_{k=1}^{N} v_k \left\{\frac{t}{k}\right\},
  \qquad t > 0,
  \label{eq:witness-wave}
\end{equation}
where $\{x\} = x - \lfloor x \rfloor$ is the fractional part. Each
term $v_k \{t/k\}$ is a sawtooth oscillation with period~$k$ and
amplitude~$v_k$. The witness wave $V(t)$ is their superposition---a
signal whose frequency content encodes the prime-weighted coefficients.

\begin{principle}[The Quadratic Form Identity]
The Gram quadratic form equals the time-domain integral of the squared
witness wave (proved in \lean{quad\_form\_time\_domain}):
\begin{equation}
  v^T G v = \int_1^{\infty} \frac{V(t)^2}{t^2}\,dt
  \label{eq:quad-form-td}
\end{equation}
This identity transforms the algebraic statement $v^T G v < 1$
(which implies RH via the Nyman--Beurling converse) into a
\emph{physical} statement: the total energy of the witness wave,
measured with the $1/t^2$ weighting (the inverse-square law), must be
less than~1.
\end{principle}

The proof proceeds by substitution: each Gram entry $G(j,k) =
\int_0^1 \{1/(jx)\}\{1/(kx)\}\,dx$ is rewritten via $t = 1/x$ as
$\int_1^\infty \{t/j\}\{t/k\}/t^2\,dt$ (proved as
\lean{substitution\_identity}), then the double sum $\sum_{j,k} v_j v_k
G(j,k)$ is recognized as $\int V(t)^2/t^2\,dt$ by bilinearity.

\subsection{The Mean--Fluctuation Decomposition}

The witness wave's squared amplitude decomposes into a mean and a
fluctuation:
\begin{equation}
  V(t)^2 = \mu_S + \bigl(V(t)^2 - \mu_S\bigr)
  \label{eq:mean-fluct}
\end{equation}
where $\mu_S = v^T R\, v + \tfrac{1}{4}(\sum_k v_k)^2$ is the
\emph{periodicMean} (proved in \lean{mu\_S\_eq}), and $R$ is the
Ramanujan matrix.

\begin{correspondence}[Vacuum Fluctuations]
The decomposition~\eqref{eq:mean-fluct} is the number-theoretic
analogue of the Casimir effect:
\begin{itemize}
  \item $\mu_S$ is the \textbf{vacuum expectation value}---the
    mean-square amplitude of the witness wave, determined by the
    Ramanujan matrix (the ``classical'' part).
  \item $V(t)^2 - \mu_S$ is the \textbf{vacuum fluctuation}---the
    deviation from the mean, which oscillates with zero time-average
    (proved in \lean{fluct\_mean\_zero}).
\end{itemize}
The Gram quadratic form splits accordingly
(proved in \lean{quad\_form\_mean\_fluct}):
\[
  v^T G v = \mu_S + \int_1^{\infty}
  \frac{V(t)^2 - \mu_S}{t^2}\,dt
\]
The first term is the classical Ramanujan energy; the second is the
quantum correction from vacuum fluctuations.
\end{correspondence}

\subsection{The Ramanujan Invariance: An Algebraic Miracle}

The deepest result in the Time-Domain Bridge is the
\emph{Ramanujan invariance} (proved in \lean{ramanujan\_invariance},
marked \textbf{THE ALGEBRAIC MIRACLE} in the code):
\begin{equation}
  R(L/j,\, L/k) = R(j, k)
  \label{eq:ramanujan-invariance}
\end{equation}
for any $L$ divisible by both $j$ and $k$, where $R(j,k)$ is the
Ramanujan matrix entry.

\begin{principle}[Hidden Scaling Symmetry]
The Ramanujan invariance is a \emph{hidden scaling symmetry} of
the Gram matrix. The periodicMean of the witness wave does not
change when we rescale the lattice by any common multiple $L$---the
``interaction strength'' between modes $j$ and $k$ depends only on
their arithmetic relationship, not on the absolute scale.

In physics, this is a \emph{conformal symmetry}: the correlation
function is invariant under dilations. The Ramanujan matrix is
conformally invariant on the integer lattice.

The proof chain composes three certified steps:
\begin{enumerate}
  \item \textbf{CoV (Change of Variables):} The integral identity
    $\int_0^L \{t/j\}\{t/k\}\,dt = \int_0^1 \{Lt/j\}\{Lt/k\}\,dt
    \cdot L$ (scaling).
  \item \textbf{Glass (Periodicity):} The integrand $\{t/j\}\{t/k\}$
    is periodic with period $\mathrm{lcm}(j,k)$ (proved in
    \lean{witnessWave\_periodic}).
  \item \textbf{Ramanujan:} The integral per period is exactly the
    Ramanujan entry $R(j,k)$ (from the Glass quadratic form identity).
\end{enumerate}
\end{principle}

\subsection{Periodicity and the IBP Identity}

The witness wave $V(t)$ is periodic with period $N!$ (proved in
\lean{witnessWave\_periodic}), as is $V(t)^2$ (proved in
\lean{witnessWave\_sq\_periodic}) and the fluctuation $V(t)^2 - \mu_S$
(proved in \lean{fluct\_periodic}).

The periodicity of the fluctuation, combined with its zero mean
(proved in \lean{fluct\_mean\_zero}), enables integration by parts
on the improper integral $\int_1^\infty (V^2 - \mu_S)/t^2\,dt$.
The IBP identity (proved in \lean{exact\_ibp\_identity}) yields:
\begin{equation}
  v^T G v = 2\mu_S - \frac{S^2}{3} + 2\int_1^\infty
  \frac{E_S(t)}{t^3}\,dt
  \label{eq:ibp-identity}
\end{equation}
where $S = \sum_k v_k$ and $E_S$ is the fluctuation primitive
(bounded by periodicity + continuity, proved in
\lean{globalFluctPrimitive\_bounded}).

\begin{correspondence}[The IBP as Virial Theorem]
The IBP identity~\eqref{eq:ibp-identity} is the number-theoretic
\emph{virial theorem}: it relates the time-averaged kinetic energy
($v^T G v$) to the potential energy ($2\mu_S - S^2/3$) plus a
correction from the boundary fluctuations. The $1/t^3$ weight in the
remainder integral ensures convergence---the fluctuations are damped
by the cube of the distance, not just the square.
\end{correspondence}
